\documentclass[11pt,a4paper]{article}
\usepackage[utf8]{inputenc}
\usepackage[spanish,es-noquoting]{babel}
\usepackage{amsmath}
\usepackage{xurl}
\usepackage{hyperref}
\usepackage{listings}
\usepackage{xcolor}
\usepackage{geometry}
\usepackage{fancyhdr}

\hypersetup{
    colorlinks=true,
    linkcolor=blue,
    urlcolor=blue,
    breaklinks=true
}

\geometry{margin=2.5cm, top=3.5cm}
\setlength{\headheight}{15pt}

\pagestyle{plain}

% Estilo de la primera pagina
\fancypagestyle{firstpage}{
  \fancyhf{}
  \fancyhead[L]{\large \textbf{Ayudantía 1 EDA}}
  \fancyhead[R]{\large \textbf{Ayudante:} Felipe Maturana.}
  \renewcommand{\headrulewidth}{0.4pt}
  \fancyfoot[C]{\thepage}
}

% Configuracion de colores para el codigo
\definecolor{codegray}{rgb}{0.5,0.5,0.5}
\definecolor{codepurple}{rgb}{0.58,0,0.82}
\definecolor{backcolour}{rgb}{0.95,0.95,0.92}

\lstdefinestyle{mystyle}{
    backgroundcolor=\color{backcolour},   
    commentstyle=\color{codegray},
    keywordstyle=\color{blue},
    numberstyle=\tiny\color{codegray},
    stringstyle=\color{codepurple},
    basicstyle=\ttfamily\footnotesize,
    breakatwhitespace=false,         
    breaklines=true,                 
    captionpos=b,                    
    keepspaces=true,                 
    numbers=none,                    
    showspaces=false,                
    showstringspaces=false,
    showtabs=false,                  
    tabsize=4
}
\lstset{style=mystyle}

\begin{document}
\thispagestyle{firstpage}

\section*{Introducción a la Complejidad Temporal}
\textbf{Lectura complementaria:} \url{https://www.geeksforgeeks.org/dsa/complete-guide-on-complexity-analysis/}

\vspace{0.2cm}
La complejidad temporal de un algoritmo se define como la cantidad de tiempo que tarda en ejecutarse en función del tamaño de la entrada ($n$). Representa el crecimiento del tiempo de ejecución respecto a la entrada, no el tiempo real en una máquina específica.

\paragraph{Notación Big O:}
\[
O(g(n)) = \{ f(n) : \exists\, c > 0, \, n_0 > 0 \text{ tales que } 0 \le f(n) \le c \cdot g(n), \; \forall n \ge n_0 \}
\]

Asumiremos que las operaciones básicas (comparaciones, asignaciones, retornos y lectura de variables) toman tiempo constante: $O(1)$. El tiempo total es la suma de los tiempos de cada declaración.

\section*{Ejemplos Básicos}

\subsection*{Ejemplo 1}
\textbf{Enunciado:} Verificar si un número $N$ es par o impar.

\textbf{Respuesta:}
\begin{lstlisting}[language=Java]
public static void checkEvenOdd(int N) {
    int r = N % 2;
    if (r == 0) {
        System.out.println("Even");
    } else {
        System.out.println("Odd");
    }
}
\end{lstlisting}

\subsection*{Ejemplo 2}
\textbf{Enunciado:} Implementar búsqueda binaria recursiva para encontrar la posición de un elemento $x$ en un arreglo ordenado. Retornar $-1$ si no se encuentra.

\textbf{Respuesta:}
\begin{lstlisting}[language=Java]
static int binarySearch(int arr[], int l, int r, int x) {
    if (r >= l) {
        int mid = l + (r - l) / 2;
        if (arr[mid] == x) return mid;
        if (arr[mid] > x) return binarySearch(arr, l, mid - 1, x);
        return binarySearch(arr, mid + 1, r, x);
    }
    return -1;
}
\end{lstlisting}

\subsection*{Ejemplo 3}
\textbf{Enunciado:} Búsqueda lineal para encontrar la posición de un elemento $x$ en un arreglo. Retornar $-1$ si no se encuentra.

\textbf{Respuesta:}
\begin{lstlisting}[language=Java]
static int search(int[] arr, int N, int x) {
    for (int i = 0; i < N; i++) {
        if (arr[i] == x) {
            return i;
        }
    }
    return -1;
}
\end{lstlisting}

\subsection*{Ejemplo 4}
\textbf{Enunciado:} Dado un arreglo $A$ y un valor $x$, determinar si existe un par de elementos en el arreglo cuya suma sea exactamente $x$.

\textbf{Respuesta:}
\begin{lstlisting}[language=Java]
static boolean chkPair(int[] A, int x) {
    int size = A.length;
    for (int i = 0; i < size - 1; i++) {
        for (int j = i + 1; j < size; j++) {
            if (A[i] + A[j] == x) {
                return true;
            }
        }
    }
    return false;
}
\end{lstlisting}

\section*{Ejercicios de Análisis}

\textbf{Ejercicio 1:} Analice el siguiente código y determine su complejidad temporal.
\begin{lstlisting}[language=Java]
for (int i = 0; i < n; i++) {
    // Operaciones O(1)
}
for (int j = 0; j < m; j++) {
    // Operaciones O(1)
}
\end{lstlisting}
\textbf{Respuesta:} La complejidad temporal es $O(n + m)$.

\vspace{0.3cm}
\textbf{Ejercicio 2:} Analice el siguiente código y determine su complejidad temporal.
\begin{lstlisting}[language=Java]
while (n > 0) {
    n = n / 2;
    // Operaciones O(1)
}
\end{lstlisting}
\textbf{Respuesta:} La complejidad temporal es $O(\log n)$.

\vspace{0.3cm}
\textbf{Ejercicio 3:} Analice el siguiente código y determine su complejidad temporal.
\begin{lstlisting}[language=Java]
for (int i = 0; i < 1000; i += 10) {
    // Operaciones O(1)
}
\end{lstlisting}
\textbf{Respuesta:} La complejidad temporal es $O(1)$, ya que el ciclo se ejecuta una cantidad constante de veces (100 iteraciones).

\newpage
\section*{Competencia}

\subsection*{Problema 1: Subtract the Product and Sum of Digits of an Integer}
\textbf{Enlace:} \url{https://leetcode.com/problems/subtract-the-product-and-sum-of-digits-of-an-integer/description/}

\textbf{Enunciado:} Dado un número entero \texttt{n}, devuelva la diferencia entre el producto de sus dígitos y la suma de sus dígitos.

\textbf{Ejemplo:}
\begin{itemize}
    \item \textbf{Entrada:} \texttt{n = 234}
    \item \textbf{Salida:} \texttt{15}
    \item \textbf{Explicación:} Producto $= 2 \times 3 \times 4 = 24$, Suma $= 2 + 3 + 4 = 9$. Resultado $= 24 - 9 = 15$.
\end{itemize}

\textbf{Respuesta:}
\begin{lstlisting}[language=Java]
class Solution {
    public int subtractProductAndSum(int n) {
        int productVal = 1;
        int sumVal = 0;
        
        while (n > 0) {
            int digit = n % 10;
            productVal *= digit;
            sumVal += digit;
            n /= 10;
        }
        
        return productVal - sumVal;
    }
}
\end{lstlisting}
\textbf{Complejidad:} Temporal $O(\log_{10} n)$, Espacial $O(1)$.

\subsection*{Problema 2: Valid Palindrome}
\textbf{Enlace:} \url{https://leetcode.com/problems/valid-palindrome/description/}

\textbf{Enunciado:} Una frase es un palíndromo si, tras convertir todo a minúsculas y eliminar los caracteres no alfanuméricos, se lee igual hacia adelante y hacia atrás. Dado un string \texttt{s}, devuelva \texttt{true} si es un palíndromo, o \texttt{false} de lo contrario.

\textbf{Ejemplo:}
\begin{itemize}
    \item \textbf{Entrada:} \texttt{s = "A man, a plan, a canal: Panama"}
    \item \textbf{Salida:} \texttt{true}
    \item \textbf{Explicación:} Eliminando caracteres no alfanuméricos queda \texttt{"amanaplanacanalpanama"}, el cual es un palíndromo.
\end{itemize}

\textbf{Respuesta:}
\begin{lstlisting}[language=Java]
class Solution {
    public boolean isPalindrome(String s) {
        int left = 0;
        int right = s.length() - 1;
        
        while (left < right) {
            while (left < right && !Character.isLetterOrDigit(s.charAt(left))) {
                left++;
            }
            while (left < right && !Character.isLetterOrDigit(s.charAt(right))) {
                right--;
            }
            
            char leftChar = Character.toLowerCase(s.charAt(left));
            char rightChar = Character.toLowerCase(s.charAt(right));
            
            if (leftChar != rightChar) {
                return false;
            }
            
            left++;
            right--;
        }
        
        return true;
    }
}
\end{lstlisting}
\textbf{Complejidad:} Temporal $O(n)$, Espacial $O(1)$.

\subsection*{Problema 3: Move Zeroes}
\textbf{Enlace:} \url{https://leetcode.com/problems/move-zeroes/description/}

\textbf{Enunciado:} Dado un arreglo de enteros \texttt{nums}, mueva todos los \texttt{0} al final mientras mantiene el orden relativo de los elementos distintos de cero. Esto debe hacerse \textit{in-place} sin realizar una copia del arreglo.

\textbf{Ejemplo:}
\begin{itemize}
    \item \textbf{Entrada:} \texttt{nums = [0,1,0,3,12]}
    \item \textbf{Salida:} \texttt{[1,3,12,0,0]}
\end{itemize}

\textbf{Respuesta:}
\begin{lstlisting}[language=Java]
class Solution {
    public void moveZeroes(int[] nums) {
        int insertPos = 0;
        
        for (int i = 0; i < nums.length; i++) {
            if (nums[i] != 0) {
                nums[insertPos] = nums[i];
                insertPos++;
            }
        }
        
        while (insertPos < nums.length) {
            nums[insertPos] = 0;
            insertPos++;
        }
    }
}
\end{lstlisting}
\textbf{Complejidad:} Temporal $O(n)$, Espacial $O(1)$.

\end{document}
